% GENERATED FILE. Edit canonical lessons, then run npm run generate:books.
% canonical-source-hash: fnv1a64:ba8a12ec0b5a1057
% canonical-lessons: JA-C248-shinjiru, JA-C248-uso, JA-C248-yakusoku, JA-C248-himitsu, JA-C248-utagau

\chapter{Promises, Lies and Secrets}
\label{ch:ja-promises-lies-and-secrets}
\hlchaptermodality{248}

\begin{chapteropening}
\textbf{By the end of this chapter:} \emph{I can ask about promises and secrets and say what someone does not believe.}

Complete the last lesson of chapter 248: I can ask about promises and secrets and say what someone does not believe.
\end{chapteropening}

\section[\texorpdfstring{\ja{しんじる} (shinjiru)}{しんじる (shinjiru)}]{\ja{しんじる} (shinjiru) — to believe; to trust}

\label{lesson:JA-C248-shinjiru}



\begin{quote}
Before the new one: say the Japanese for to be in trouble, then the Japanese for courage.
\end{quote}

\subsection*{You'll want to know: \ja{しんじる}}
\textbf{\ja{しんじる}} — \emph{shinjiru} — \textquotedblleft{}to believe; to trust\textquotedblright{}.

\begin{quote}
\textbf{\ja{かれの} \ja{こたえを} \ja{しんじますか。}} — \emph{kare no kotae o shinjimasu ka} — \textquotedblleft{}Do you believe his answer?\textquotedblright{}
\end{quote}

Polite \textbf{\ja{しんじます}}, \emph{shinjimasu}; an ichidan verb.

\begin{cousinweb}
\textbf{\ja{しんじる}} (\emph{shinjiru}) is a mixed word, built on Sino-Japanese \emph{shin}, \textquotedblleft{}trust\textquotedblright{} (Middle Chinese \emph{sinH}), borrowed from Chinese, with a native verb ending \emph{-jiru} (JMdict, KANJIDIC2; Middle Chinese in Baxter's transcription).
\end{cousinweb}

\begin{grammarlens}[title={What you've built}]
One more word for everyday life at A2.
\end{grammarlens}

\subsection*{Your turn}
\begin{itemize}
\raggedright
  \item \emph{Say it:} \emph{shinjiru}
  \item \emph{Say it:} \emph{shinjiru}, once more
  \item \emph{Say it:} \emph{kare no kotae o shinjimasu ka} — Do you believe his answer?
  \item \emph{Recall:} say the Japanese for to be in trouble, then the Japanese for courage, then say \emph{shinjiru} again
\end{itemize}

\begin{tcolorbox}[breakable,colback=teal!4,colframe=teal!35!black,title={Before you move on}]
What does \ja{しんじる} mean? (\textquotedblleft{}To believe; to trust\textquotedblright{}.) Say it once more.
\end{tcolorbox}
\section[\texorpdfstring{\ja{うそ} (uso)}{うそ (uso)}]{\ja{うそ} (uso) — a lie; a fib}

\label{lesson:JA-C248-uso}



\begin{quote}
Before the new one: say the Japanese for courage, then the Japanese for to believe.
\end{quote}

\subsection*{You'll want to know: \ja{うそ}}
\textbf{\ja{うそ}} — \emph{uso} — \textquotedblleft{}a lie; a fib\textquotedblright{}.

\begin{quote}
\textbf{\ja{かれは} \ja{うそを} \ja{いいません。}} — \emph{kare wa uso o iimasen} — \textquotedblleft{}He does not tell lies.\textquotedblright{}
\end{quote}

\begin{cousinweb}
\textbf{\ja{うそ}} (\emph{uso}) is a native Japanese word (wago), inherited rather than borrowed from Chinese; its kanji was attached later as a kun'yomi, and its older origin is not settled (JMdict, KANJIDIC2).
\end{cousinweb}

\begin{grammarlens}[title={What you've built}]
One more word for everyday life at A2.
\end{grammarlens}

\subsection*{Your turn}
\begin{itemize}
\raggedright
  \item \emph{Say it:} \emph{uso}
  \item \emph{Say it:} \emph{uso}, once more
  \item \emph{Say it:} \emph{kare wa uso o iimasen} — He does not tell lies.
  \item \emph{Recall:} say the Japanese for courage, then the Japanese for to believe, then say \emph{uso} again
\end{itemize}

\begin{tcolorbox}[breakable,colback=teal!4,colframe=teal!35!black,title={Before you move on}]
What does \ja{うそ} mean? (\textquotedblleft{}A lie; a fib\textquotedblright{}.) Say it once more.
\end{tcolorbox}
\section[\texorpdfstring{\ja{やくそく} (yakusoku)}{やくそく (yakusoku)}]{\ja{やくそく} (yakusoku) — a promise; an appointment}

\label{lesson:JA-C248-yakusoku}



\begin{quote}
Before the new one: say the Japanese for to believe, then the Japanese for a lie.
\end{quote}

\subsection*{You'll want to know: \ja{やくそく}}
\textbf{\ja{やくそく}} — \emph{yakusoku} — \textquotedblleft{}a promise; an appointment\textquotedblright{}.

\begin{quote}
\textbf{\ja{あしたの} \ja{やくそくは} \ja{ごごですか。}} — \emph{ashita no yakusoku wa gogo desu ka} — \textquotedblleft{}Is tomorrow's appointment in the afternoon?\textquotedblright{}
\end{quote}

\begin{cousinweb}
\textbf{\ja{やくそく}} (\emph{yakusoku}) is Sino-Japanese (kango): two Chinese characters read with their on'yomi, readings borrowed from Chinese. They are \emph{yaku}, \textquotedblleft{}promise\textquotedblright{} (Middle Chinese \emph{'jak}) and \emph{soku}, \textquotedblleft{}bundle, tie\textquotedblright{} (\emph{syowk}) (KANJIDIC2; Middle Chinese in Baxter's transcription).
\end{cousinweb}

\begin{grammarlens}[title={What you've built}]
One more word for everyday life at A2.
\end{grammarlens}

\subsection*{Your turn}
\begin{itemize}
\raggedright
  \item \emph{Say it:} \emph{yakusoku}
  \item \emph{Say it:} \emph{yakusoku}, once more
  \item \emph{Say it:} \emph{ashita no yakusoku wa gogo desu ka} — Is tomorrow's appointment in the afternoon?
  \item \emph{Recall:} say the Japanese for to believe, then the Japanese for a lie, then say \emph{yakusoku} again
\end{itemize}

\begin{tcolorbox}[breakable,colback=teal!4,colframe=teal!35!black,title={Before you move on}]
What does \ja{やくそく} mean? (\textquotedblleft{}A promise; an appointment\textquotedblright{}.) Say it once more.
\end{tcolorbox}
\section[\texorpdfstring{\ja{ひみつ} (himitsu)}{ひみつ (himitsu)}]{\ja{ひみつ} (himitsu) — a secret}

\label{lesson:JA-C248-himitsu}



\begin{quote}
Before the new one: say the Japanese for a lie, then the Japanese for a promise.
\end{quote}

\subsection*{You'll want to know: \ja{ひみつ}}
\textbf{\ja{ひみつ}} — \emph{himitsu} — \textquotedblleft{}a secret\textquotedblright{}.

\begin{quote}
\textbf{\ja{これは} \ja{かぞくの} \ja{ひみつですか。}} — \emph{kore wa kazoku no himitsu desu ka} — \textquotedblleft{}Is this a family secret?\textquotedblright{}
\end{quote}

\begin{cousinweb}
\textbf{\ja{ひみつ}} (\emph{himitsu}) is Sino-Japanese (kango): two Chinese characters read with their on'yomi, readings borrowed from Chinese. They are \emph{hi}, \textquotedblleft{}hidden\textquotedblright{} (Middle Chinese \emph{pijH}) and \emph{mitsu}, \textquotedblleft{}close, dense\textquotedblright{} (\emph{mit}) (KANJIDIC2; Middle Chinese in Baxter's transcription).
\end{cousinweb}

\begin{grammarlens}[title={What you've built}]
One more word for everyday life at A2.
\end{grammarlens}

\subsection*{Your turn}
\begin{itemize}
\raggedright
  \item \emph{Say it:} \emph{himitsu}
  \item \emph{Say it:} \emph{himitsu}, once more
  \item \emph{Say it:} \emph{kore wa kazoku no himitsu desu ka} — Is this a family secret?
  \item \emph{Recall:} say the Japanese for a lie, then the Japanese for a promise, then say \emph{himitsu} again
\end{itemize}

\begin{tcolorbox}[breakable,colback=teal!4,colframe=teal!35!black,title={Before you move on}]
What does \ja{ひみつ} mean? (\textquotedblleft{}A secret\textquotedblright{}.) Say it once more.
\end{tcolorbox}
\section[\texorpdfstring{\ja{うたがう} (utagau)}{うたがう (utagau)}]{\ja{うたがう} (utagau) — to doubt; to suspect}

\label{lesson:JA-C248-utagau}



\begin{quote}
Before the new one: say the Japanese for a promise, then the Japanese for a secret.
\end{quote}

\subsection*{You'll want to know: \ja{うたがう}}
\textbf{\ja{うたがう}} — \emph{utagau} — \textquotedblleft{}to doubt; to suspect\textquotedblright{}.

\begin{quote}
\textbf{\ja{ははは} \ja{ちちの} \ja{こたえを} \ja{うたがいません。}} — \emph{haha wa chichi no kotae o utagaimasen} — \textquotedblleft{}My mother does not doubt my father's answer.\textquotedblright{}
\end{quote}

Polite \textbf{\ja{うたがいます}}, \emph{utagaimasu}; a godan verb.

\begin{cousinweb}
\textbf{\ja{うたがう}} (\emph{utagau}) is a native Japanese word (wago), inherited rather than borrowed from Chinese; its kanji was attached later as a kun'yomi (JMdict, KANJIDIC2).
\end{cousinweb}

\begin{grammarlens}[title={What you've built}]
One more word for everyday life at A2.
\end{grammarlens}

\subsection*{Your turn}
\begin{itemize}
\raggedright
  \item \emph{Say it:} \emph{utagau}
  \item \emph{Say it:} \emph{utagau}, once more
  \item \emph{Say it:} \emph{haha wa chichi no kotae o utagaimasen} — My mother does not doubt my father's answer.
  \item \emph{Recall:} say the Japanese for a promise, then the Japanese for a secret, then say \emph{utagau} again
\end{itemize}

\begin{tcolorbox}[breakable,colback=teal!4,colframe=teal!35!black,title={Before you move on}]
What does \ja{うたがう} mean? (\textquotedblleft{}To doubt; to suspect\textquotedblright{}.) Say it once more.
\end{tcolorbox}
